\begin{description}
\item[(A)] \textbf{Transiting Exoplanet.} The figure below presents the light curve of an exoplanet transit. On the X-axis are represented the days elapsed since the first observation, and on the Y-axis is shown the percentage of the average luminosity by which the brightness of the system increases or decreases relative to the average brightness of the star, that is,
\[
100\times\frac{L(t)-L_{\mathrm{average}}}{L_{\mathrm{average}}}.
\]
Determine the period $P$ of the planet and the ratio Planet Radius/Star Radius.
\ioaafig{2025_TH_Q09-f1.pdf}
On the Y-axis, 0.2 means 0.2\% and 1.0 means 1\%.
\item[(B)] \textbf{Eclipsing Binary Star System.} Another type of object with an interesting light curve is an eclipsing binary star system. The lightcurve of such an object is presented below, with the X-axis representing days since the start of the observing mission, and the Y-axis representing the relative brightness of the system, that is, $L(t)/L_{\mathrm{average}}$. We know that the smaller star is hotter than the larger star. Calculate the period $P$ of the system and the ratio Small Star Radius/Big Star Radius.
\ioaafig{2025_TH_Q09-f2.pdf}
On the Y-axis, 0.7 means 70\% and 1.0 means 100\%.
\end{description}
